\section{Polarization and exact prime--pole interaction}
\label{sec:interaction}

For \(y>1/2\), the two translates have disjoint support.  Hermitian
polarization, with the first form variable antilinear, gives
\begin{align*}
\qform[h_y^+]
&=\frac12\{\qform[h_y^{\mathrm R}]+\qform[h_y^{\mathrm L}]\}
+\Re\qform(h_y^{\mathrm R},h_y^{\mathrm L}),\\
\qform[h_y^-]
&=\frac12\{\qform[h_y^{\mathrm R}]+\qform[h_y^{\mathrm L}]\}
-\Re\qform(h_y^{\mathrm R},h_y^{\mathrm L}).
\end{align*}
Inversion symmetry makes the cross term real.  Thus
\begin{equation}
\boxed{\Delta_h(y)=2\qform(h_y^{\mathrm R},h_y^{\mathrm L}).}
\label{eq:polarization}
\end{equation}

Let
\[
M_h=\widehat h(i/2)=\widehat h(-i/2)
=\int_\R h(x)e^{x/2}\,dx>0.
\]
The two pole terms in \eqref{eq:preform} contribute exactly
\begin{equation}
M_h^2(e^y+e^{-y}).
\label{eq:pole-cross}
\end{equation}
The prime-power cross term is
\begin{equation}
-
\sum_{p,m\ge1}
\frac{\log p}{p^{m/2}}
K_h(2y-m\log p).
\label{eq:prime-cross}
\end{equation}
Because \(K_h\) is supported on \([-1,1]\), only
\[
2y-1<m\log p<2y+1
\]
contributes.  Equality is zero because \(K_h(\pm1)=0\).  This is a support
fact, not an endpoint half-weight convention.

Define the full Chebyshev measure
\[
d\Psi=\sum_{n\ge2}\Lambda(n)\delta_n,\qquad
\psi(X)=\sum_{n\le X}\Lambda(n),
\]
and
\[
dE=d\Psi-dx.
\]
This is the prime-power error distribution associated with
\(\psi(X)-X\), not the primes-only distribution associated with
\(\theta(X)-X\).

The continuous term can be evaluated without approximation:
\begin{align}
\int_0^\infty
x^{-1/2}K_h(2y-\log x)\,dx
&=
e^y\int_{-1}^{1}K_h(r)e^{-r/2}\,dr
\notag\\
&=e^yM_h^2.
\label{eq:continuous-main}
\end{align}
Therefore
\begin{align}
\Ih(y)
&=
e^yM_h^2
-
\sum_{n\ge2}
\frac{\Lambda(n)}{\sqrt n}
K_h(2y-\log n),
\label{eq:I-physical}
\end{align}
and the exact cross interaction is
\begin{equation}
\qform(h_y^{\mathrm R},h_y^{\mathrm L})
=
A_h^\infty(y)+e^{-y}M_h^2+\Ih(y),
\label{eq:Q-decomposition}
\end{equation}
where
\[
A_h^\infty(y)
=
\frac{1}{2\pi}
\int_\R m_\infty(t)e^{2iyt}|\widehat h(t)|^2\,dt.
\]
Equation~\eqref{eq:Q-decomposition} is the exact prime--pole cancellation:
the \(e^yM_h^2\) pole main term has been combined with the continuous part
of \(dE\) before any estimate is made.

For a packet realized inside
\([-\log\lambda,\log\lambda]\), the geometric condition is
\[
y+\frac12\le\log\lambda.
\]
Every contributing prime power then automatically satisfies the inclusive
semilocal cutoff \(p^m\le\lambda^2\).
